\chapter{Calibrazione dei parametri e simulazione}\label{chap:simulazione}

Nei capitoli precedenti abbiamo costruito il modello e scelto una forma analitica per ciascuna delle sue funzioni. Per poterlo simulare, tuttavia, è necessario assegnare un valore a tutti i parametri. In questo capitolo scegliamo tali valori e descriviamo la dinamica che ne deriva. \par

Non disponendo di dati reali con cui confrontare il modello, procediamo nel modo seguente: fissiamo alcuni requisiti che la dinamica deve rispettare, ad esempio quanto può calare il gruppo o quanti atleti possono essere infortunati, e cerchiamo i valori dei parametri che li soddisfano. In altre parole, invece di domandarci quali siano i valori ``veri'' dei parametri, che non possiamo conoscere senza dati, ci chiediamo quali valori producano un comportamento del gruppo coerente con quello che vogliamo rappresentare. \par

Lo scenario che costruiamo è sfavorevole, ma non estremo: rappresenta un periodo di allenamento intenso senza alcuna forma di prevenzione, con molti infortuni, in cui però il gruppo resta in gran parte unito e gli atleti sani raramente abbandonano. Nel capitolo successivo verrà introdotto un controllo, cioè le strategie di prevenzione a disposizione dell'allenatore, e questo scenario servirà come termine di confronto per misurarne l'efficacia. Se lo scenario di partenza presentasse pochi infortuni, il controllo avrebbe poco margine per migliorarlo e sarebbe difficile confrontare strategie diverse. Per questo motivo preferiamo partire da una situazione in cui gli infortuni sono un problema evidente.

\section{Impostazione della simulazione}

Fissiamo un orizzonte di $T = 50$ sedute di allenamento. Ipotizzando una seduta al giorno, tale orizzonte equivale a circa sette settimane, vale a dire un periodo di preparazione di durata media, paragonabile alla preparazione per una gara importante oppure a un raduno collegiale. La durata scelta è sufficientemente ampia da permettere di osservare l'accumulo della fatica e degli infortuni, ma al tempo stesso abbastanza contenuta da rispettare l'ipotesi secondo cui gli atleti infortunati non hanno il tempo di guarire del tutto. \par

Le condizioni iniziali sono:
\begin{gather*}
    S(0) = 400 \\
    P(0) = 100 \\
    I(0) = M(0) = 0 \\
    \bar{F}_f(0) = \bar{F}_p(0) = 0
\end{gather*}
Il gruppo di partenza è dunque formato da $N(0) = 500$ atleti, dei quali l'$80\%$ sani e il $20\%$ con infortuni pregressi. Inizialmente nessun atleta risulta infortunato e il gruppo è riposato, poiché assumiamo che il periodo cominci dopo una fase di recupero. La numerosità del gruppo è tale da consentire di trattare le popolazioni come numeri reali, come richiesto dalle ipotesi del modello. \par

La capienza del tatami è $A_{\max} = 800$ atleti. All'inizio del periodo la densità sul tatami è quindi pari a $0.625$: il tatami è affollato, ma non saturo. Nel corso del periodo la densità cala, poiché una parte degli atleti riporta infortuni gravi e smette di allenarsi e, nelle ultime settimane, alcuni abbandonano: alla fine essa scende a circa $0.48$.

\subsection{Carico di allenamento}

Il carico di allenamento costituisce una scelta dell'allenatore e non viene calibrato. Adottiamo il carico periodico presentato nel capitolo sulle funzioni, nel quale si susseguono una seduta ad alto carico, una a carico medio e una priva di carico:
\[
    c_f(t) = c_p(t) = \begin{cases}
        1 & \text{se } t = 3n \\
        0.5 & \text{se } t = 3n + 1 \\
        0 & \text{se } t = 3n + 2
    \end{cases}
\]
Tale carico corrisponde a un ciclo di tre giorni assai diffuso nella pratica: un giorno di randori intenso, un giorno di tecnica e una seduta leggera, di scarico, con esercizi a basso impegno come le cadute o lo stretching. Il carico medio per seduta risulta pari a $0.5$. \par

\begin{figure}[ht]
    \centering
    \includegraphics[width = 0.5\textwidth]{pictures/grip2.jpg}
    \caption[Esercizio a coppie]{Due judoka si preparano a un esercizio a coppie durante una seduta di tecnica, in cui il carico è intermedio.}
\end{figure}

Dato che carico fisico e carico psicologico coincidono, anche le due componenti della fatica sono uguali per ogni $t$, ovvero $\bar{F}_f(t) = \bar{F}_p(t)$. Questa scelta rende più semplice la lettura dei risultati e non compromette gli scopi del capitolo. La distinzione tra le due componenti assumerà rilevanza nel momento in cui si metteranno a confronto programmi di allenamento caratterizzati da un diverso equilibrio tra impegno fisico e impegno mentale. \par

È opportuno osservare che i parametri individuati valgono per questo specifico carico. Qualora l'allenatore modificasse il programma, i parametri del gruppo rimarrebbero invariati, ma la dinamica cambierebbe: proprio questo consentirà, in seguito, di confrontare programmi di allenamento differenti.

\section{Parametri delle funzioni}

I parametri si articolano in due gruppi. Il primo comprende i parametri delle funzioni $\mu$, $\nu$, $\lambda$ e $r$, che ne definiscono la forma. Il secondo comprende i parametri delle equazioni del modello, ovvero $\alpha$, $\pi$, $\rho$, i moltiplicatori di abbandono $\beta_\cdot$ e i moltiplicatori di rischio $\gamma_\cdot$. \par

I due gruppi svolgono un ruolo differente. I parametri delle funzioni descrivono \emph{in che modo} variano nel tempo ingressi, abbandoni, rischio e recupero. I parametri del modello descrivono invece \emph{quanto} tali fenomeni incidono su ciascuna popolazione. Per tale ragione li trattiamo separatamente: i parametri delle funzioni vengono fissati direttamente, in funzione del comportamento che intendiamo ottenere, mentre i parametri del modello sono calcolati nella sezione seguente. La separazione dei due gruppi consente inoltre di ridurre il numero di parametri da calcolare, da diciannove a dieci. \par

I valori scelti per i parametri delle funzioni sono riportati nella \autoref{tab:sim-funzioni} e commentati qui di seguito.

\begin{table}[ht]
    \centering
    \caption{Parametri delle funzioni del modello}\label{tab:sim-funzioni}
    \begin{tblr}{
        colspec = {cccl},
        hline{4,5,9} = {0.03em},
    }
        Funzione & Parametro & Valore & Significato \\
        \SetCell[r=2]{c} $\mu$ & $\omega_\mu$ & $5$ & ingressi nella prima seduta \\
              & $\sigma_\mu$ & $0.15$ & durata del periodo di ingresso \\
        $\nu$ & $\sigma_\nu$ & $0.15$ & durata del periodo di abbandono \\
        \SetCell[r=4]{c}$\lambda$ & $\omega_f$ & $0.5$ & peso della fatica fisica \\
                  & $\omega_p$ & $0.3$ & peso della fatica psicologica \\
                  & $\omega_d$ & $0.2$ & peso della densità \\
                  & $b$ & $1$ & soglia \\
        \SetCell[r=2]{c}$r$ & $\omega_r$ & $1$ & recupero iniziale \\
            & $\sigma_r$ & $0.05$ & velocità di decadimento del recupero \\
    \end{tblr}
\end{table}

\paragraph{Funzione di crescita}

Nel corso della prima seduta entrano $5$ nuovi atleti, dopodiché gli ingressi calano rapidamente. Complessivamente entrano circa $49$ nuovi atleti, ovvero circa il $10\%$ del gruppo iniziale, e quasi tutti, all'incirca $47$, arrivano entro le prime $15$ sedute. Superata la metà del periodo gli ingressi risultano praticamente nulli. \par

Tale profilo descrive un gruppo che si completa nelle prime due settimane: alcuni atleti giungono in ritardo, ad esempio per ragioni di lavoro, di studio o di trasporto, ma dopo questa fase il gruppo rimane chiuso. La crescita della popolazione nella fase iniziale del periodo dipende quasi esclusivamente da questi due parametri, poiché nelle prime settimane gli abbandoni sono trascurabili.

\paragraph{Funzione di abbandono}

Gli abbandoni presentano un andamento opposto. Il rapporto tra la funzione di abbandono e il suo valore massimo è praticamente nullo a metà del periodo, assume valore circa $0.08$ dopo due terzi del periodo, circa $0.41$ all'$80\%$ e circa $0.80$ al $90\%$, fino a raggiungere il massimo nell'ultima seduta. Gli abbandoni si concentrano dunque nelle ultime due settimane. \par

Ciò rispecchia quanto spesso si verifica nella pratica: verso la conclusione di un periodo di allenamento prolungato crescono gli impegni esterni, la stanchezza accumulata si fa sentire e alcuni atleti scelgono di interrompere anzitempo. L'entità con cui la funzione incide su ciascuna popolazione dipende dai moltiplicatori $\beta_\cdot$, calcolati nella sezione successiva.

\paragraph{Funzione di infortunio}

Il fattore di rischio più rilevante è la fatica fisica, seguita dalla fatica psicologica e dalla densità sul tatami. I pesi $\omega_f$, $\omega_p$ e $\omega_d$ rispecchiano esattamente questo ordine. La soglia $b = 1$ garantisce che, con il gruppo riposato, il rischio sia contenuto ma non nullo. \par

Per valutare l'entità del rischio è utile considerare il rapporto $\frac{\lambda}{\gamma_\cdot}$, che denominiamo rischio relativo: esso indica la frazione del rischio massimo $\gamma_\cdot$ effettivamente raggiunta ed è identico per tutte le popolazioni. Nella simulazione il rischio relativo cresce parallelamente alla fatica, passando da circa $0.29$ all'inizio a circa $0.79$ al termine del periodo. \par

Il rischio relativo si mantiene pertanto distante sia da zero sia dal suo valore massimo, pari a $1$. Ciò è rilevante: in questa fascia il rischio risulta sensibile alle variazioni della fatica, e dunque una modifica del carico di allenamento si traduce in una variazione effettiva del numero di infortuni.

\paragraph{Funzione di recupero}

La scelta più significativa concerne la velocità $\sigma_r$ con cui il recupero diminuisce. Nel capitolo sulle funzioni avevamo impiegato, a titolo di esempio, $\sigma_r = 0.2$. Con tale valore, tuttavia, il recupero si riduce quasi a zero già a metà del periodo: gli atleti non sono più in grado di smaltire la fatica, che cresce ininterrottamente fino a circa $18$. Dopo circa $25$ sedute il rischio relativo raggiunge $0.98$, cioè praticamente il massimo, come illustrato nella \autoref{fig:sim-fatica} e nella \autoref{fig:sim-rischio}. \par

Da quel momento il rischio non dipende più dalla fatica, e quindi neppure dal carico di allenamento: qualsiasi programma di allenamento condurrebbe al medesimo risultato. Con $\sigma_r = 0.2$, inoltre, gli infortuni salgono a circa $709$, rispetto ai circa $428$ ottenuti con il valore scelto più avanti: un incremento di due terzi imputabile esclusivamente al recupero. \par

\begin{figure}[ht]
    \centering
    \includegraphics[width = \textwidth]{simulazione/fatica.pdf}
    \caption{Fatica media del gruppo per due valori di $\sigma_r$}\label{fig:sim-fatica}
\end{figure}

Scegliamo pertanto $\sigma_r = 0.05$. Il recupero cala più lentamente: vale $1$ all'inizio, circa $0.76$ dopo $10$ sedute, circa $0.45$ dopo $25$ e ancora $0.15$ al termine del periodo. La fatica oscilla seguendo il ciclo di tre sedute: cresce sempre dopo il giorno di randori e diminuisce sempre dopo la seduta leggera. Dopo il giorno di tecnica, invece, per gran parte del periodo diminuisce, poiché il recupero smaltisce più fatica di quanta ne aggiunga il carico medio; soltanto nelle ultime sedute, dalla trentasettesima in poi, quando il recupero è più basso, aumenta anche dopo la tecnica. Nelle prime venti sedute non supera circa $1.2$, poi cresce gradualmente, man mano che il recupero diminuisce, fino a circa $2.8$ al termine del periodo. \par

\begin{figure}[ht]
    \centering
    \includegraphics[width = \textwidth]{simulazione/rischio.pdf}
    \caption{Rischio relativo di infortunio per due valori di $\sigma_r$}\label{fig:sim-rischio}
\end{figure}

Con questa scelta il rischio relativo segue lo stesso andamento oscillante della fatica e si mantiene lontano dal massimo. In tal modo il carico di allenamento continua a incidere sugli infortuni per l'intera durata del periodo, che è esattamente ciò che intendiamo studiare quando introdurremo il controllo.

\section{Parametri del modello}

Rimangono da determinare i dieci parametri del modello: $\alpha$, $\pi$, $\rho$, $\beta_S$, $\beta_P$, $\beta_M$, $\beta_I$, $\gamma_S$, $\gamma_P$ e $\gamma_M$. Li scegliamo in modo tale che la simulazione soddisfi i seguenti requisiti:
\begin{itemize}
    \item la popolazione complessiva non superi il $10\%$ di crescita:
        \[
            N(t) \leq 1.10 \cdot N(0)
        \]
    \item al termine del periodo la popolazione complessiva risulti pari all'$85\%$ di quella iniziale:
        \[
            N(T) = 0.85 \cdot N(0)
        \]
    \item in ciascuna seduta gli infortunati e i malconci non superino la metà del gruppo:
        \[
            I(t) + M(t) \leq 0.50 \cdot N(t)
        \]
    \item gli infortunati gravi non eccedano il $10\%$ del gruppo:
        \[
            I(t) / N(t) \leq 0.1
        \]
    \item i moltiplicatori rispettino l'ordine $\beta_S \leq \beta_P \leq \beta_M \leq \beta_I$ e $\gamma_S \leq \gamma_P \leq \gamma_M$.
\end{itemize}

Ogni requisito possiede un significato ben preciso. Il primo stabilisce che il gruppo è già quasi completo all'inizio del periodo e che i nuovi arrivi sono esigui. Il secondo determina l'entità della riduzione: alla conclusione del periodo se ne va circa un atleta su sette, una perdita consistente ma non tale da compromettere il gruppo. Il terzo esclude le situazioni in cui la maggioranza del gruppo risulta infortunata e l'allenamento perderebbe di significato. Il quarto definisce la gravità dello scenario: nel momento peggiore un atleta su dieci è fermo a causa di un infortunio grave. L'ultimo riprende le ipotesi del modello: chi ha un infortunio pregresso o in corso abbandona con maggiore facilità e si infortuna più facilmente di chi è sano. \par

I requisiti relativi al calo della popolazione e agli infortunati gravi sono fissati esattamente al loro valore limite, affinché lo scenario risulti il più sfavorevole possibile senza oltrepassare tali limiti. Gli altri requisiti costituiscono limiti massimi: vanno rispettati, ma non è necessario raggiungerli.

\subsection{Valori a priori}

Per ciascun parametro individuiamo un valore iniziale, denominato valore a priori, e un intervallo di valori ammessi, riportati nella \autoref{tab:sim-intervalli}. Gli intervalli escludono i valori estremi, quali un rischio nullo o un recupero immediato, che non descriverebbero una situazione realistica. Per i sani e i pregressi gli intervalli dei moltiplicatori di abbandono sono molto ristretti: un atleta che sta bene e si allena con regolarità difficilmente lascia il gruppo nel giro di poche settimane.

\begin{table}[ht]
    \centering
    \caption{Intervalli e valori a priori dei parametri del modello}\label{tab:sim-intervalli}
    \begin{tblr}{
        colspec = {ccc},
    }
        Parametro & Intervallo & Valore a priori \\
        $\alpha$ & $[0.5, 1]$ & $0.8$ \\
        $\pi$ & $[0.1, 0.6]$ & $0.4$ \\
        $\rho$ & $[0.02, 0.3]$ & $0.05$ \\
        $\beta_S$ & $[0.001, 0.02]$ & $0.005$ \\
        $\beta_P$ & $[0.001, 0.03]$ & $0.01$ \\
        $\beta_M$ & $[0.005, 0.4]$ & $0.05$ \\
        $\beta_I$ & $[0.005, 0.4]$ & $0.1$ \\
        $\gamma_S$ & $[0.001, 0.3]$ & $0.02$ \\
        $\gamma_P$ & $[0.001, 0.3]$ & $0.04$ \\
        $\gamma_M$ & $[0.001, 0.3]$ & $0.08$ \\
    \end{tblr}
\end{table}

I valori a priori descrivono un gruppo che si allena in assenza di prevenzione:
\begin{itemize}
    \item $\alpha = 0.8$: i nuovi atleti presentano la stessa composizione del gruppo iniziale, con l'$80\%$ di sani;
    \item $\pi = 0.4$: due infortuni su cinque sono gravi;
    \item $\rho = 0.05$: gli infortunati gravi tornano ad allenarsi lentamente, dopo circa $20$ sedute in media;
    \item i moltiplicatori di abbandono sono molto bassi per i sani e per i pregressi, $\beta_S = 0.005$ e $\beta_P = 0.01$, e più elevati per i malconci e gli infortunati, $\beta_M = 0.05$ e $\beta_I = 0.1$: chi sta bene rimane quasi sempre nel gruppo, chi si è fatto male è molto più incline a lasciarlo;
    \item i moltiplicatori di rischio aumentano dai sani ai malconci: un malconcio presenta un rischio quattro volte superiore a quello di un sano.
\end{itemize}

Con i valori a priori la simulazione è già prossima ai requisiti, ma non li rispetta completamente. Il picco di infortunati gravi è pari a circa il $10.2\%$ del gruppo, appena al di sopra del limite del $10\%$, e al termine del periodo resta l'$87\%$ del gruppo, cioè $436$ atleti, anziché l'$85\%$. Gli infortuni totali sono circa $417$. Occorre pertanto correggere leggermente i valori a priori.

\subsection{Calcolo dei parametri}

Per correggere i valori a priori abbiamo impiegato la funzione \texttt{minimize} della libreria Python SciPy \cite{virtanen2020}, una delle librerie più diffuse per il calcolo scientifico. Tale funzione ricerca i valori dei parametri che minimizzano una determinata quantità, denominata funzione obiettivo, rispettando al contempo una serie di vincoli. \par

Come funzione obiettivo abbiamo utilizzato la distanza dai valori a priori. Per ogni parametro si calcola la differenza tra il valore proposto e il valore a priori, la si divide per la larghezza dell'intervallo del parametro, la si eleva al quadrato e infine si sommano i risultati:
\[
    \text{distanza} = \sum_{\text{parametri}} \left(\frac{\text{valore} - \text{valore a priori}}{\text{larghezza dell'intervallo}}\right)^2.
\]
La divisione per la larghezza dell'intervallo rende confrontabili parametri che possono variare in intervalli di ampiezza diversa: ogni variazione viene misurata in proporzione allo spazio a disposizione di quel parametro. Il quadrato, dal canto suo, fa sì che le differenze positive e negative contribuiscano allo stesso modo e che una correzione grande su un solo parametro pesi più di molte correzioni piccole. \par

I requisiti sono stati passati alla funzione come vincoli: i due requisiti fissati esattamente, ovvero il calo della popolazione e il picco di infortunati gravi, come vincoli di uguaglianza, gli altri come vincoli di disuguaglianza. Gli intervalli della \autoref{tab:sim-intervalli} sono stati passati come limiti per ciascun parametro. In presenza di vincoli, la funzione \texttt{minimize} adopera automaticamente il metodo SLSQP\footnote{Il metodo SLSQP (Sequential Least Squares Programming, o Sequential Least Squares Quadratic Programming) è un algoritmo di ottimizzazione numerica progettato per minimizzare funzioni scalari di più variabili soggette a vincoli non lineari.}, adatto a problemi con pochi parametri e vincoli di questo tipo. \par

Il risultato consiste nei parametri più vicini ai valori a priori tra tutti quelli che rispettano i requisiti. In tal modo la stima iniziale viene corretta il meno possibile, e i parametri finali conservano il significato attribuito ai valori a priori. La procedura è riassunta nella \autoref{fig:calibrazione}.

\begin{figure}[ht]
    \centering
    \begin{tikzpicture}[
        block/.style = {
            rectangle, rounded corners = 4pt, draw = ColorP, line width = 1pt, fill = ColorP!12,
            minimum width = 2.5cm, minimum height = 1.4cm, align = center, font = \small
        },
        input/.style = {
            rectangle, rounded corners = 4pt, draw = gray, line width = 0.8pt, fill = gray!8,
            minimum width = 2.5cm, minimum height = 1cm, align = center, font = \small
        },
        flow/.style = {-Stealth, line width = 1.2pt}
    ]
        \node[input] (prior) {Valori a priori\\e intervalli};
        \node[block, right = 0.6cm of prior] (minimize) {\texttt{minimize}\\distanza minima\\dai valori\\a priori};
        \node[block, draw = ColorS, fill = ColorS!12, right = 0.6cm of minimize] (calibrated) {Parametri\\calibrati};
        \node[block, right = 0.6cm of calibrated] (simulation) {Simulazione\\finale};

        \node[input, above = 0.7cm of minimize] (requirements) {Requisiti (vincoli)};
        \node[input, below = 0.9cm of minimize] (trial) {Simulazione di prova};

        \draw[flow] (prior) -- (minimize);
        \draw[flow] (minimize) -- (calibrated);
        \draw[flow] (calibrated) -- (simulation);
        \draw[flow, draw = gray] (requirements) -- (minimize);
        \draw[flow, draw = gray] ([xshift = -0.5cm] minimize.south) to[bend right = 35]
            node[left, font = \footnotesize] {parametri proposti} ([xshift = -0.5cm] trial.north);
        \draw[flow, draw = gray] ([xshift = 0.5cm] trial.north) to[bend right = 35]
            node[right, font = \footnotesize] {requisiti rispettati?} ([xshift = 0.5cm] minimize.south);
    \end{tikzpicture}
    \caption{La procedura di calibrazione dei parametri}\label{fig:calibrazione}
\end{figure}



I parametri calcolati sono riportati nella \autoref{tab:sim-calibrati}. Essi risultano tutti vicini ai valori a priori: la maggior parte cambia di pochi millesimi. \par

I moltiplicatori di abbandono dei sani e dei pregressi restano praticamente invariati. Le correzioni interessano soprattutto i malconci e gli infortunati. Il moltiplicatore di abbandono dei malconci $\beta_M$ passa da $0.050$ a $0.059$ e quello degli infortunati $\beta_I$ da $0.100$ a $0.105$: abbandonano un po' di più, e ciò porta la popolazione finale dall'$87\%$ all'$85\%$ richiesto. Il tasso di recupero parziale $\rho$ passa da $0.050$ a $0.054$: gli infortunati gravi tornano ad allenarsi un po' prima, in media dopo circa $18.7$ sedute anziché $20$, e questo abbassa il picco di infortunati gravi al $10\%$. \par

Con questi parametri, nel momento di massima fatica un atleta sano ha una probabilità di infortunarsi in una singola seduta di circa l'$1.7\%$, un atleta con infortunio pregresso di circa il $3.2\%$ e un malconcio di circa il $6.2\%$.

\begin{table}[ht]
    \centering
    \caption{Parametri del modello: valori a priori e valori calibrati}\label{tab:sim-calibrati}
    \begin{tblr}{
        colspec = {ccc},
    }
        Parametro & Valore a priori & Valore calibrato \\
        $\alpha$ & $0.800$ & $0.800$ \\
        $\pi$ & $0.400$ & $0.397$ \\
        $\rho$ & $0.050$ & $0.054$ \\
        $\beta_S$ & $0.005$ & $0.005$ \\
        $\beta_P$ & $0.010$ & $0.010$ \\
        $\beta_M$ & $0.050$ & $0.059$ \\
        $\beta_I$ & $0.100$ & $0.105$ \\
        $\gamma_S$ & $0.020$ & $0.021$ \\
        $\gamma_P$ & $0.040$ & $0.040$ \\
        $\gamma_M$ & $0.080$ & $0.078$ \\
    \end{tblr}
\end{table}

La \autoref{tab:sim-confronto} confronta i risultati ottenuti con i valori a priori e con i valori calibrati. Le correzioni sono modeste e lo scenario muta poco: il gruppo finale passa da $436$ a $425$ atleti, il picco di infortunati gravi resta intorno ai $52$ atleti e gli infortuni totali passano da circa $417$ a circa $428$. \\
Con infortuni totali non si intendono le singole persone che si sono infortunate, bensì il numero complessivo di infortuni verificatisi. Infatti lo stesso atleta può infortunarsi più di una volta.

\begin{table}[ht]
    \centering
    \caption{Risultati della simulazione con i valori a priori e con i valori calibrati}\label{tab:sim-confronto}
    \begin{tblr}{
        colspec = {lcc},
    }
        & Valori a priori & Valori calibrati \\
        Popolazione massima & $110\%$ & $110\%$ \\
        Popolazione finale & $87\%$ & $85\%$ \\
        Picco di infortunati gravi & $53$ ($10.2\%$) & $52$ ($10\%$) \\
        Picco di malconci & $126$ & $131$ \\
        Massimo di infortunati e malconci & $35\%$ & $36\%$ \\
        Infortuni totali & $417$ & $428$ \\
    \end{tblr}
\end{table}

\subsection{Popolazione totale}

La \autoref{fig:sim-totale} mostra la popolazione totale rispetto a quella iniziale. Si possono distinguere tre fasi. \par

Nella prima fase, nel corso delle prime due settimane, il gruppo cresce grazie ai nuovi ingressi: dopo $10$ sedute è già al $108\%$ di quello iniziale. Nella seconda fase, nella parte centrale del periodo, il gruppo si mantiene stabile intorno al $110\%$: gli ingressi sono terminati e gli abbandoni non sono ancora iniziati. Il massimo, circa $549$ atleti, viene raggiunto intorno alla venticinquesima seduta. \par

Nella terza fase, nell'ultimo terzo del periodo, iniziano gli abbandoni e il gruppo cala. Intorno alla quarantaquattresima seduta torna sotto il valore iniziale, alla quarantacinquesima è al $97\%$ e alla fine arriva a $425$ atleti, cioè l'$85\%$ di quello iniziale, come richiesto. \par

\begin{figure}[ht]
    \centering
    \includegraphics[width = \textwidth]{simulazione/popolazione_totale.pdf}
    \caption{Popolazione totale rispetto a quella iniziale}\label{fig:sim-totale}
\end{figure}

Complessivamente abbandonano circa $124$ atleti: soltanto $12$ sani e $4$ pregressi, contro $44$ infortunati e $64$ malconci. Gli infortunati e i malconci, che all'inizio non esistevano, sono dunque responsabili dell'$87\%$ degli abbandoni. Gli infortuni costituiscono la causa principale della riduzione del gruppo: ridurre gli infortuni significherebbe anche mantenere unito il gruppo.

\subsection{Le quattro popolazioni}

La \autoref{fig:sim-popolazioni} mostra le quattro popolazioni e la \autoref{tab:sim-popolazioni} ne riporta i valori ogni dieci sedute.

\begin{figure}[ht]
    \centering
    \includegraphics[width = \textwidth]{simulazione/popolazioni.pdf}
    \caption{Le quattro popolazioni con i parametri calibrati}\label{fig:sim-popolazioni}
\end{figure}

\begin{table}[ht]
    \centering
    \caption{Le quattro popolazioni ogni dieci sedute}\label{tab:sim-popolazioni}
    \begin{tblr}{
        colspec = {cccccc},
    }
        Seduta & $S$ & $P$ & $I$ & $M$ & Infortunati e malconci \\
        $0$ & $400$ & $100$ & $0$ & $0$ & $0\%$ \\
        $10$ & $399$ & $93$ & $17$ & $32$ & $9\%$ \\
        $20$ & $370$ & $79$ & $31$ & $68$ & $18\%$ \\
        $30$ & $333$ & $65$ & $44$ & $106$ & $27\%$ \\
        $40$ & $293$ & $51$ & $52$ & $131$ & $35\%$ \\
        $50$ & $243$ & $36$ & $38$ & $108$ & $34\%$ \\
    \end{tblr}
\end{table}

Nelle prime sedute i sani rimangono quasi costanti, poiché i nuovi ingressi compensano gli infortuni; successivamente, sia i sani sia i pregressi diminuiscono per tutto il periodo. La diminuzione è dovuta soprattutto agli infortuni, che li spostano tra i malconci o tra gli infortunati, mentre gli abbandoni incidono poco. I pregressi diminuiscono più rapidamente dei sani in proporzione, perché presentano un rischio di infortunio più elevato: alla fine ne restano $36$ su $100$, mentre i sani sono $243$ su $400$. \par

Gli infortunati e i malconci partono da zero e crescono per gran parte del periodo. Raggiungono il massimo intorno alla quarantesima seduta, con circa $52$ infortunati e $131$ malconci, per poi diminuire a causa degli abbandoni. I malconci superano i pregressi già intorno alla ventitreesima seduta, ma infortunati e malconci insieme restano sempre meno numerosi di sani e pregressi insieme. \par

I malconci sono sempre più numerosi degli infortunati gravi, per due ragioni. La prima è che la maggior parte degli infortuni è lieve e conduce l'atleta direttamente in $M$. La seconda è che anche gli infortunati gravi, quando guariscono parzialmente, passano in $M$. I malconci, a loro volta, presentano il rischio di infortunio più alto e possono infortunarsi gravemente, tornando in $I$. \par

Al termine del periodo restano $243$ sani, $36$ pregressi, $38$ infortunati e $108$ malconci: circa un terzo degli atleti rimasti è infortunato o malconcio.

\subsection{Infortunati}

La \autoref{fig:sim-frazioni} mostra la frazione di infortunati gravi e quella di infortunati e malconci insieme. Entrambe crescono in modo regolare per gran parte del periodo. Gli infortunati gravi raggiungono il $10\%$ del gruppo intorno alla quarantaduesima seduta, senza mai superarlo. Infortunati e malconci insieme arrivano al massimo al $36\%$ del gruppo, rimanendo quindi ben al di sotto del limite del $50\%$. Nelle ultime sedute entrambe le frazioni diminuiscono leggermente, non perché vi siano meno infortuni, ma perché molti infortunati e malconci abbandonano. \par

\begin{figure}[ht]
    \centering
    \includegraphics[width = \textwidth]{simulazione/frazione_infortunati.pdf}
    \caption{Frazione di infortunati gravi e di infortunati e malconci}\label{fig:sim-frazioni}
\end{figure}

Complessivamente si verificano circa $428$ infortuni, di cui $170$ gravi e $258$ lievi. In media si verificano circa $9$ infortuni per seduta, ma il numero cresce nel tempo insieme alla fatica: sono circa $4$ nella prima seduta e arrivano a circa $13$ intorno alla quarantatreesima. Nella seconda metà del periodo si verificano circa $270$ infortuni, contro i $158$ della prima metà. \par

Gli infortuni si distribuiscono tra le popolazioni nel modo seguente: circa $184$ riguardano i sani, $70$ i pregressi e $174$ i malconci. I malconci, pur essendo molto meno numerosi dei sani, producono il $41\%$ degli infortuni, quasi quanto i sani, perché costituiscono la popolazione più a rischio. La \autoref{fig:sim-infortuni-popolazione} mostra questa ripartizione. \par

\begin{figure}[ht]
    \centering
    \includegraphics[width = \textwidth]{simulazione/infortuni_per_popolazione.pdf}
    \caption{Infortuni per popolazione di provenienza}\label{fig:sim-infortuni-popolazione}
\end{figure}


Una parte degli atleti si infortuna più di una volta, ad esempio prima in modo lieve, diventando malconci, e poi in modo grave. Rapportando gli infortuni al numero di presenze sul tatami si ottengono circa $17$ infortuni ogni $1000$ presenze, un valore elevato, coerente con uno scenario privo di prevenzione.

\section{Analisi di sensibilità}\label{sec:sensibilita}

I parametri del modello non sono stimati a partire da dati reali, ma scelti con la calibrazione per plausibilità. È quindi importante capire quanto i risultati dipendano dai valori scelti. Per farlo abbiamo ripetuto la simulazione variando un parametro alla volta del $20\%$ in più e in meno rispetto al suo valore, lasciando invariati tutti gli altri. Questo approccio, detto ``one-at-a-time'', è il metodo più semplice di analisi di sensibilità~\cite{pianosi2016}. Abbiamo considerato quattordici parametri: i moltiplicatori di rischio e di abbandono dei malconci e degli infortunati, la probabilità $\pi$, il tasso di recupero parziale $\rho$, i parametri della funzione di infortunio, la capienza del tatami, la velocità di decadimento del recupero e l'ampiezza degli ingressi. \par

La \autoref{fig:sim-sensibilita} mostra la variazione degli infortuni totali, con i parametri ordinati dal più influente al meno influente, e la \autoref{tab:sim-sensibilita} riporta i valori.

\begin{figure}[ht]
    \centering
    \includegraphics[width = \textwidth]{simulazione/sensibilita.pdf}
    \caption{Variazione degli infortuni totali al variare di ciascun parametro del $20\%$}\label{fig:sim-sensibilita}
\end{figure}

\begin{table}[ht]
    \centering
    \caption{Infortuni totali e riduzione con la prevenzione al variare di ciascun parametro del $20\%$}\label{tab:sim-sensibilita}
    \begin{tblr}{
        colspec = {X[c, m] X[c, m] X[c, m] X[c, m] X[c, m]},
        width = \textwidth,
        hline{2} = {0pt},
        hline{2} = {2-3}{leftpos = -1, rightpos = -1, endpos, 0.05em},
        hline{2} = {4-5}{leftpos = -1, rightpos = -1, endpos, 0.05em},
        hline{3} = {0.05em},
    }
        \SetCell[r=2]{c} Parametro & \SetCell[c=2]{c} Infortuni totali & & \SetCell[c=2]{c} Riduzione con prevenzione & \\
        & $-20\%$ & $+20\%$ & $-20\%$ & $+20\%$ \\
        $\gamma_S$ & $377$ & $474$ & $31.2\%$ & $30.4\%$ \\
        $b$ & $471$ & $385$ & $31.0\%$ & $30.7\%$ \\
        $\sigma_r$ & $391$ & $469$ & $30.7\%$ & $31.0\%$ \\
        $\gamma_M$ & $397$ & $457$ & $30.0\%$ & $31.5\%$ \\
        $\omega_f$ & $403$ & $452$ & $31.0\%$ & $30.7\%$ \\
        $\gamma_P$ & $411$ & $442$ & $31.2\%$ & $30.5\%$ \\
        $\omega_p$ & $413$ & $443$ & $31.0\%$ & $30.8\%$ \\
        $\pi$ & $442$ & $415$ & $31.4\%$ & $30.3\%$ \\
        $\omega_\mu$ & $421$ & $435$ & $30.9\%$ & $30.8\%$ \\
        $\rho$ & $421$ & $434$ & $30.5\%$ & $31.1\%$ \\
        $A_{\max}$ & $435$ & $424$ & $30.9\%$ & $30.9\%$ \\
        $\omega_d$ & $423$ & $433$ & $30.9\%$ & $30.9\%$ \\
        $\beta_M$ & $431$ & $425$ & $30.9\%$ & $30.8\%$ \\
        $\beta_I$ & $428$ & $428$ & $30.9\%$ & $30.8\%$ \\
    \end{tblr}
\end{table}

Nessuna variazione modifica gli infortuni totali di più del $12\%$: con parametri variati del $20\%$, gli infortuni restano compresi tra circa $377$ e $474$, contro i $428$ dello scenario di riferimento. Il modello, quindi, non amplifica le incertezze sui parametri. \par

I parametri più influenti sono quelli del rischio e del recupero. Il primo è il moltiplicatore di rischio dei sani $\gamma_S$, perché i sani sono la popolazione più numerosa. Seguono la soglia $b$, che con un valore più alto riduce il rischio di base, la velocità di decadimento del recupero $\sigma_r$, il moltiplicatore di rischio dei malconci $\gamma_M$ e il peso della fatica fisica $\omega_f$. I parametri degli abbandoni, $\beta_M$ e $\beta_I$, la capienza del tatami $A_{\max}$ e il peso della densità $\omega_d$ incidono invece per meno del $2\%$: gli abbandoni si concentrano nelle ultime sedute, quando la maggior parte degli infortuni è già avvenuta, e la densità varia poco nel corso del periodo. \par

Un risultato meno intuitivo riguarda la probabilità $\pi$ che un infortunio sia grave: aumentandola, gli infortuni totali diminuiscono. Con più infortuni gravi, infatti, più atleti si fermano in $I$, dove non sono esposti al rischio, e meno atleti diventano malconci, che sono la popolazione più a rischio. Il numero di infortuni gravi aumenta, ma il numero complessivo diminuisce. \par

Anche la popolazione finale varia poco: in tutti i casi resta compresa tra l'$83\%$ e l'$87\%$ di quella iniziale. \par

Infine, l'analisi riguarda anche la prevenzione che introdurremo nel capitolo successivo. La riduzione degli infortuni ottenuta con la prevenzione è praticamente insensibile ai parametri: in tutti i casi resta compresa tra il $30.0\%$ e il $31.5\%$.

\section{Considerazioni finali}

Lo scenario ottenuto descrive un periodo di allenamento intenso senza prevenzione: si verificano circa $428$ infortuni, il gruppo perde il $15\%$ degli atleti e alla fine circa un terzo degli atleti rimasti è infortunato o malconcio. Al tempo stesso, lo scenario rispetta tutti i requisiti: la popolazione non cresce più del $10\%$, alla fine è l'$85\%$ di quella iniziale, gli infortunati gravi non superano il $10\%$ e gli infortunati e i malconci non superano la metà del gruppo. \par

Dai risultati emergono tre aspetti sui quali il controllo potrà intervenire:
\begin{itemize}
    \item il rischio di infortunio cresce con la fatica, e difatti gli infortuni sono più numerosi nella seconda metà del periodo: ridurre la fatica, ad esempio modificando l'alternanza tra sedute intense e sedute di recupero, riduce gli infortuni;
    \item i malconci sono la popolazione più a rischio e producono circa il $41\%$ degli infortuni: una strategia di prevenzione mirata a loro, ad esempio riducendone il carico, può avere un effetto rilevante;
    \item gli infortunati e i malconci sono la causa principale degli abbandoni: ridurre gli infortuni aiuta anche a mantenere unito il gruppo.
\end{itemize}

I parametri scelti non sono stimati a partire da dati reali, bensì dai requisiti che abbiamo fissato e dai valori a priori. Con requisiti o valori a priori differenti si otterrebbe uno scenario diverso. Lo scenario va dunque inteso come un caso di riferimento, utile per confrontare strategie diverse, e non come una previsione di ciò che accade in una palestra reale. L'analisi di sensibilità mostra comunque che variazioni moderate dei parametri non cambiano la forma della dinamica né le conclusioni qualitative. Disporre anche di pochi dati reali, come il numero di iscritti, di abbandoni e di infortuni in una stagione, permetterebbe di scegliere i requisiti in modo più oggettivo.
